\documentclass[11pt,a4paper]{article}
\usepackage[utf8]{inputenc}
\usepackage[T1]{fontenc}
\usepackage[french]{babel}
\usepackage{geometry}
\geometry{margin=2.2cm}
\usepackage{enumitem}
\usepackage{xcolor}
\usepackage{titlesec}
\usepackage{booktabs}
\usepackage{amssymb}
\usepackage{hyperref}
\hypersetup{colorlinks=true, linkcolor=blue!60!black, urlcolor=blue!60!black}
\usepackage{tcolorbox}

\definecolor{pekgno}{RGB}{22,101,52}
\definecolor{todo}{RGB}{150,60,0}
\definecolor{codebg}{RGB}{245,245,245}

\titleformat{\section}{\Large\bfseries\color{pekgno}}{\thesection}{0.6em}{}
\titleformat{\subsection}{\large\bfseries\color{black}}{\thesubsection}{0.6em}{}

\newlist{tasks}{enumerate}{1}
\setlist[tasks,1]{label=\textcolor{todo}{$\square$}~, leftmargin=1.6em, itemsep=4pt}
\newlist{details}{enumerate}{1}
\setlist[details,1]{label=\textbullet, leftmargin=1.6em, itemsep=2pt}

\newcommand{\taskitem}[1]{\item #1}
\newcommand{\codefile}[1]{\texttt{\detokenize{#1}}}
\newcommand{\code}[1]{{\small\color{black!80}\texttt{#1}}}

\begin{document}

\begin{center}
{\LARGE\bfseries\color{pekgno} PEKEGNO --- Liste des tâches de développement}\\[0.4em]
{\large Évolutions Phase 4 : Employés, Tranches de paiement, Primes, Comptabilité, Abonnements, Séminaire, Bilan du jour}\\[0.3em]
{\small Document de travail à destination des développeurs --- Spécifications détaillées : \codefile{SPECIFICATION_EVOLUTIONS.md}}
\end{center}

\vspace{0.8em}
\noindent\textbf{Légende des priorités :} \textcolor{pekgno}{\textbf{P1}} = urgent/critique (bug), \textcolor{todo}{\textbf{P2}} = socle métier, \textcolor{black}{\textbf{P3}} = fonctionnalité autonome.

\begin{tcolorbox}[colback=codebg, colframe=pekgno, title=Rappels transverses]
\begin{details}
\item Backend Laravel dans \codefile{backend/}, frontend React 19 + TS dans \codefile{frontend/src/}.
\item Toute chaîne visible passe par i18n (\codefile{fr.ts} et \codefile{en.ts}).
\item Toute action métier a sa permission ; les nouvelles permissions sont seedées dans \codefile{backend/database/seeders/RoleSeeder.php}.
\item Les exports actuels sont CSV (\codefile{backend/app/Http/Controllers/Api/ExportController.php}) ; les nouveaux exports devront produire de vrais fichiers \codefile{.xlsx} (PhpSpreadsheet / maatwebsite/excel).
\item Le menu est centralisé dans \codefile{frontend/src/components/layout/navItems.ts} ; le sous-menu d'agence dans \codefile{frontend/src/components/agencies/AgencyLayout.tsx} ; les routes dans \codefile{frontend/src/router/index.tsx}.
\end{details}
\end{tcolorbox}

\section{Correction : recherche de factures (P1)}

\begin{tasks}
\taskitem Exécuter la migration manquante \codefile{2026_08_14_000001_add_client_name_to_invoices.php} (colonne \code{client\_name}) sur la base : \code{php artisan migrate}; vérifier \code{migrate:status}.
\taskitem Renforcer la recherche dans \codefile{backend/app/Http/Controllers/Api/InvoiceController.php} (\code{index()}, lignes 50--54) : ajouter à la clause \code{orWhere} une recherche sur la relation \code{client} (prénom, nom, email) pour ne plus dépendre de la seule colonne \code{client\_name}.
\taskitem Ajouter le \textbf{filtre par client} sur la liste des factures \codefile{frontend/src/pages/invoices/InvoiceListPage.tsx} : un \code{Select} « Client » (champ \code{client\_id} déjà transmis à l'API, ligne 78) alimenté par \code{clientsApi.list()}.
\taskitem Ajouter la clé i18n \code{invoices.filterClient} (fr/en).
\end{tasks}

\section{Menu Employés (P3)}

\subsection{Conception}
Un employé (ex. caissier) peut vendre et être commissionné comme un commercial, mais n'en est pas un. \textbf{Option retenue} : colonne discriminante \code{kind} sur \code{commercials} (\code{'commercial'} | \code{'employe'}) pour réutiliser toute la logique primes/points.

\begin{tasks}
\taskitem Migration : \codefile{2026_08_15_000002_add_kind_to_commercials_table.php} --- \code{enum('kind', ['commercial','employe'])->default('commercial')} après \code{agency\_id}.
\taskitem Backend \codefile{app/Models/Commercial.php} : \code{kind} dans \code{\$fillable} + \code{scopeKind()}.
\taskitem Backend \codefile{app/Http/Controllers/Api/CommercialController.php} : filtre \code{kind} dans \code{index()} (et dans \code{search}, \code{ranking}, \code{stats} si nécessaire).
\taskitem Routes \codefile{backend/routes/api.php} : alias \code{/employees} vers les mêmes contrôleurs (index, search, available-users, ranking, stats, points, store, show, update, destroy), permissions \code{employes.*}.
\taskitem Seeder \codefile{RoleSeeder.php} : permissions \code{employes.consulter/creer/modifier/supprimer/exporter} pour super-admin, direction, responsable-agence, comptable.
\taskitem Export \codefile{ExportController.php} : méthode \code{employees()} (gabarit \code{commercials()}) + colonne \code{kind}.
\taskitem Frontend types \codefile{src/types/commercial.ts} : \code{kind?} sur \code{Commercial}, \code{CommercialListParams}, \code{CommercialPayload}.
\taskitem Frontend \codefile{src/api/commercials.api.ts} : \code{list} accepte \code{kind}.
\taskitem Frontend pages : rendre \codefile{CommercialListPage.tsx} et \codefile{CommercialDetailPage.tsx} paramétrables par \code{kind} (pattern \code{fixedAgencyId}).
\taskitem Routes frontend : \code{/employees}, \code{/employees/:id}, \code{/agencies/:agencyId/employees[:/:employeeId]} (wrappers pattern \codefile{AgencyCommercialsPage.tsx}).
\taskitem Nav : entrée \code{nav.employees} dans \codefile{navItems.ts} + sous-menu agence dans \codefile{AgencyLayout.tsx}.
\taskitem i18n : bloc \code{employees.*} (fr/en).
\end{tasks}

\section{Paiement en tranches (max 3) et historique des versements (P2)}

\begin{tasks}
\taskitem Backend \codefile{InvoiceController::pay()} (lignes 245--294) : interdire un 4\ieme{} paiement --- \textbf{décision actée : simple plafond de 3 paiements par facture, avance comprise} (\code{payments()->count() >= 3} $\Rightarrow$ erreur 422 « Paiement en tranches limité à 3 »). Faire de même dans \code{applyPayment()} (lignes 321--344).
\taskitem Frontend : tableau \textbf{Historique des versements} dans \codefile{frontend/src/pages/invoices/InvoiceDetailPage.tsx} à partir de \code{invoice.payments} (montant, méthode, avance?, \textbf{date}, encaissé par, commentaire). Réutiliser les clés i18n existantes \code{invoices.paymentHistory}, \code{paymentAmount}, \code{paymentMethod}, \code{paymentDate}, \code{paymentReceivedBy}, \code{paymentIsAdvance}.
\taskitem Frontend modal « Encaisser » : message i18n si 3 paiements atteints ; permettre la saisie de \code{paid\_at}.
\end{tasks}

\section{Primes à l'encaissement + prime fixe par service (P2)}

\subsection{Règle métier}
Chaque encaissement (tranche) génère la commission du commercial/employé sur le \textbf{montant encaissé} (pas au paiement intégral). Si un service a une \textbf{prime fixe}, elle prime sur le pourcentage.

\begin{tasks}
\taskitem Migration \codefile{2026_08_15_000003_add_bonus_fixed_to_services_table.php} : \code{decimal('bonus\_fixed',12,2)->nullable()} sur \code{services}.
\taskitem Migration \codefile{2026_08_15_000004_create_commission_payments_table.php} : table \code{commission\_payments} (commercial\_id, invoice\_id, payment\_id nullable, service\_id nullable, amount, base\_amount, rule [percent/fixed/service\_fixed], rate nullable, invoice\_total, created\_by).
\taskitem Modèle \codefile{app/Models/CommissionPayment.php} + relations sur \code{Commercial}, \code{Invoice}, \code{InvoicePayment}.
\taskitem Service \codefile{app/Services/CommissionService.php} : \code{calculateForPayment(Invoice, float paid)} et \code{recordForPayment(...)} (écriture journal + log). Résolution multi-lignes : part du paiement par ligne selon son poids (\code{line\_total / invoice\_total}) ; ligne avec \code{service\_id} et \code{bonus\_fixed} non nul $\rightarrow$ calculée \textbf{à part} : \code{bonus\_fixed $\times$ paid/total} ; sinon règle du commercial sur la part de la ligne (\code{percent} : \code{rate\% $\times$ part} ; \code{fixed} : \code{value $\times$ part/line\_total}).
\taskitem \codefile{InvoiceController::pay()} (lignes 273--283) et \code{applyPayment()} (lignes 335--343) : remplacer le bloc « si paid $\rightarrow$ commission sur total » par \code{CommissionService->recordForPayment()} \textbf{à chaque encaissement} (idempotent par \code{payment\_id}). \textbf{Décisions actées} : \code{commission\_amount} conservé comme agrégat (somme des \code{commission\_payments}) ; points \textbf{inchangés} --- \code{PointsService::awardForSale} au soldé (idempotent).
\taskitem \codefile{Invoice::show} : charger \code{commissionPayments}.
\taskitem Frontend : champ « Prime fixe » dans \codefile{ServiceFormModal.tsx} + badge sur les cartes services.
\taskitem Frontend : section « Commissions versées » dans \codefile{CommercialDetailPage.tsx} (stats par tranche).
\taskitem Types + i18n : \code{Service.bonus\_fixed?}, \code{CommissionPayment}, clés \code{services.bonusFixed}, \code{commercials.commissionsHistory}.
\end{tasks}

\section{Comptabilité (P2)}

\subsection{Table des transactions}
Colonnes : N\code{o} (auto), Agence, Rubrique (nom ou \textbf{n\code{o} de facture}), Montant, Client, Date, Type (entrée/sortie), Opérateur, Objet (commentaire), Bénéficiaire (sorties), Justification (sorties).

\begin{tasks}
\taskitem Migration \codefile{2026_08_15_000005_create_accounting_tables.php} : \code{accounting\_categories} (name, type [income/expense], agency\_id nullable, is\_system) et \code{accounting\_transactions} (number auto-increment, agency\_id, category\_id, type, label, reference, amount, client\_id, invoice\_id, transacted\_at, operator\_id, note, beneficiary, justification).
\taskitem Modèles \codefile{app/Models/AccountingTransaction.php}, \codefile{AccountingCategory.php} + relations.
\taskitem Journalisation automatique : dans \codefile{InvoiceController::pay()} et \code{applyPayment()}, après chaque \code{invoice\_payments}, créer une transaction \code{income} (label \code{"Facture \{number\} --- versement"}, reference = n\code{o} facture, client, montant, \code{transacted\_at = paid\_at}, opérateur = receveur).
\taskitem Contrôleur \codefile{app/Http/Controllers/Api/AccountingController.php} : \code{index} (filtres agency\_id, type, category\_id, client\_id, from, to, search ; totaux entrées/sorties/solde), \code{store} (manuelle ; sortie $\Rightarrow$ beneficiary + justification requis), \code{update}, \code{destroy} (transactions auto non supprimables), \code{categories} CRUD.
\taskitem Routes \codefile{backend/routes/api.php} : \code{accounting/transactions}, \code{accounting/categories}, \code{/exports/accounting}. Permissions \code{comptabilite.*} + \code{accounting-categories.*}.
\taskitem Export Excel \codefile{ExportController::accounting()} (filtres identiques à index).
\taskitem Seeder : catégories système par défaut (ex. « Encaissement facture »).
\taskitem Frontend types \codefile{src/types/accounting.ts} + API \codefile{src/api/accounting.api.ts}.
\taskitem Pages : \codefile{frontend/src/pages/accounting/AccountingPage.tsx} (global, filtres + totaux + export + modales CRUD + gestion catégories) et wrapper agence \code{AgencyAccountingPage}.
\taskitem Routes : \code{/accounting} et \code{/agencies/:agencyId/accounting} ; nav global \code{nav.accounting} + sous-menu agence ; \code{ExportKind} += \code{'accounting'}.
\taskitem i18n : bloc \code{accounting.*} (fr/en).
\end{tasks}

\section{Reporting commerciaux (P3)}

\begin{tasks}
\taskitem Backend : endpoint agrégé (dans \codefile{StatsController.php} ou \codefile{CommercialReportController.php}) : filtres agency\_id, commercial\_id, from, to ; sortie par commercial/employé : ventes, CA facturé/encaissé, nb tranches, commissions par tranche, points, prospects, conversion ; totaux + classement.
\taskitem Export Excel : \code{ExportController::commercialReport()} (mêmes filtres).
\taskitem Permissions \code{commerciaux.reporting} : super-admin, direction, responsable-agence, comptable.
\taskitem Frontend : page \codefile{frontend/src/pages/commercials/CommercialReportPage.tsx} (filtres, résumé, détail par vente/tranche, points, prospects, export).
\taskitem Route \code{/commercials/report} + lien depuis \codefile{CommercialListPage.tsx}.
\taskitem i18n : bloc \code{commercials.report.*}.
\end{tasks}

\section{Abonnements (P3)}

\begin{tasks}
\taskitem Migrations \codefile{2026_08_15_000006_create_subscription_tables.php} : \code{subscription\_packs} (agency\_id, name, description, is\_active), \code{subscription\_pack\_services} (pack, service, price\_per\_month --- \textbf{4 services par pack}), \code{subscriptions} (pack, agency, client, months, price\_per\_month, total\_price, start\_date, end\_date, invoice\_id nullable).
\taskitem Modèles : \codefile{SubscriptionPack}, \codefile{SubscriptionPackService}, \codefile{Subscription}.
\taskitem Contrôleur \codefile{app/Http/Controllers/Api/SubscriptionController.php} : CRUD packs ; \code{store} abonnement $\rightarrow$ \textbf{génère une facture} (réutiliser \code{InvoiceNumberGenerator}) ; \code{renew} $\rightarrow$ \textbf{nouvelle facture}.
\taskitem Routes + permissions \code{abonnements.*}.
\taskitem Journalisation comptable : encaissement d'abonnement $\rightarrow$ entrée.
\taskitem Frontend : types \codefile{src/types/subscription.ts}, API, page \codefile{AgencySubscriptionsPage.tsx} (packs, création abonnement, renouvellement, lien vers facture).
\taskitem Routes \code{/agencies/:agencyId/subscriptions} ; sous-menu agence + menu global admin ; i18n \code{subscriptions.*}.
\end{tasks}

\section{Catégorie Séminaire --- Pass Classique / Premium / VIP (P3)}

\begin{tasks}
\taskitem Catégorie « Séminaire » seedée (ou création catalogue) ; \code{is\_seminar} sur le service.
\taskitem Migration \codefile{2026_08_15_000007_create_seminar_tiers_table.php} : \code{seminar\_tiers} (service\_id, tier [classique/premium/vip], label, price, description, unique(service\_id, tier)).
\taskitem Règle « intelligente » : prix dérivés du prix de base à la création (Classique = prix, Premium = $\times$1.5, VIP = $\times$2.5), modifiables dans le formulaire service.
\taskitem Migration \codefile{invoice\_items} : \code{string('pass\_tier')->nullable()}.
\taskitem Backend : relation \code{seminarTiers()} et \code{is\_seminar} sur \codefile{Service.php} ; CRUD tiers dans \codefile{ServiceController.php} ; \codefile{InvoiceController::store()} accepte \code{pass\_tier} par ligne et résout le prix.
\taskitem Frontend : champs prix des passes dans \codefile{ServiceFormModal.tsx} ; \code{Select} « Pass » dans \codefile{InvoiceFormPage.tsx} et \codefile{QuickSalePage.tsx}.
\taskitem Types : \code{InvoiceLinePayload.pass\_tier?}, \code{Service.seminar\_tiers} ; i18n \code{services.seminar.*}, \code{invoices.passLabel}.
\end{tasks}

\section{Bilan du jour (P3)}

\begin{tasks}
\taskitem Backend : endpoint \code{StatsController::dailyBilan} (ou \codefile{BilanController}) : par service $\rightarrow$ nombre + somme (\code{invoice\_items} sur factures non annulées du jour, ventes séminaire par pass) ; total services ; encaissements cash / mobile (\code{invoice\_payments.payment\_method} du jour) ; dépense du jour (sorties \code{accounting\_transactions}) ; solde initial (encaissé du jour précédent) ; \textbf{solde final = total encaissé $-$ dépense du jour} ; contrôle \code{cash + mobile == total services vendus}.
\taskitem Option stockage : table \code{daily\_balances} (agency\_id, date, solde\_initial, solde\_final) pour chaîner les jours.
\taskitem Export Excel : \code{ExportController::dailyBilan()}.
\taskitem Permissions \code{bilans.consulter / exporter}.
\taskitem Frontend : page \codefile{src/pages/reports/DailyBilanPage.tsx} (date + agence, tableau par service, totaux, cash/mobile, solde initial, dépense, solde final, export).
\taskitem Routes \code{/bilans} et \code{/agencies/:agencyId/bilans} ; nav global + sous-menu agence ; i18n \code{bilan.*}.
\end{tasks}

\section{Décisions actées (réponses aux questions ouvertes)}

\begin{details}
\item \textbf{Tranches} : pas de plan de tranches obligatoire --- simple plafond de 3 paiements par facture, avance comprise.
\item \textbf{Primes multi-lignes} : prime à chaque versement ; lignes à \code{bonus\_fixed} calculées à part (prime fixe proratée), autres lignes avec la règle du commercial.
\item \textbf{Points commerciaux} : inchangés --- attribués au paiement intégral (au soldé), pas par tranche.
\item \code{commission\_amount} sur \code{invoices} : \textbf{conservé comme agrégat} (somme des \code{commission\_payments}) pour compatibilité des statistiques.
\item \textbf{Factures annulées} : aucune écriture comptable générée.
\end{details}

\section{Ordre de mise en œuvre et tests}

\begin{tasks}
\taskitem P1 --- Bug recherche factures (\S1).
\taskitem P2 --- Tranches (\S3) + Primes par encaissement et prime fixe (\S4).
\taskitem P2 --- Comptabilité (\S5).
\taskitem P3 --- Bilan du jour (\S9) puis Employés (\S2), Reporting (\S6), Abonnements (\S7), Séminaire (\S8).
\end{tasks}

\textbf{Tests backend} (\codefile{backend/tests/}) : payer en 3 tranches max (refus au 4\ieme{} paiement, avance comprise), commission par tranche (percent / fixed / service\_fixed), recherche factures (numéro, client\_name, client lié), journalisation comptable auto (paiement $\rightarrow$ entrée), abonnement $\rightarrow$ facture + renouvellement, bilan du jour (cohérence cash + mobile = total).

\textbf{Tests frontend} : parcours manuels des nouvelles pages, ouverture des fichiers Excel sans corruption, i18n fr/en.

\end{document}
